% from ega0-21-fr.tex
\oldpage[0\textsubscript{IV-1}]{181}
  on a
  \[
    d(E/L_0,k/k_0)=d(E/K_0,k/k_0)=d(E/M_0,k/k_0),
  \]
  d'où, par une dernière application de (21.6.4.2),
  \[
    \delta=d(E/M_0,N/k_0),
  \]
  ce qui prouve (21.9.6.1). En outre, lorsque $k$ est remplacé par une
  sous-extension $k'$ de $k_0$, $N$ est remplacé par une sous-extension $N'$,
  donc $d(E/M_0,N'/k_0)\leqslant d(E/M_0,N/k_0)$ (21.6.4.2), ce qui démontre
  la dernière assertion de (i).

  \item[(ii)] Pour tout $\alpha$, posons
  $N_\alpha=k_\alpha((T_1^{p^2},\ldots,T_r^{p^2}))$. En vertu de (21.9.6.1)
  et de (21.8.3), il s'agit de montrer (compte tenu de ce que
  $K_0^p\subset k_\alpha$ et de ce que $E$ est une extension de type fini de
  $M_0$) que l'on a
  \[
    \bigcap_\alpha N_\alpha=k_0(M_0^p).
  \]
  Or, on a
  \begin{equation}
  \tag{21.9.6.6}
  \label{0.21.9.6.6-fr}
    \bigcap_\alpha N_\alpha=N'
    =(k_0(K_0^p))((T_1^{p^2},\ldots,T_r^{p^2}))
  \end{equation}
  en vertu de (21.8.8.3) et de la relation
  $\bigcap_\alpha k_\alpha=k_0(K_0^p)$ (21.8.6). D'autre part, on a
  $k_0(M_0^p)=k_0(K_0^p((T_1^{p^2},\ldots,T_r^{p^2})))$. Pour achever la
  démonstration de (21.9.6, (ii)), il reste donc à prouver le

  \begin{lemma}[21.9.6.7]
  \label{0.21.9.6.7-fr}
  Soient $k_0$ un corps de caractéristique $p>0$ tel que
  $[k_0:k_0^p]<+\infty$, $K_0$ une extension de $k_0$. Alors on a
  \begin{equation}
  \tag{21.9.6.8}
  \label{0.21.9.6.8-fr}
    (k_0(K_0^p))((T_1,\ldots,T_r))
    =k_0\bigl(K_0^p((T_1,\ldots,T_r))\bigr).
  \end{equation}
  \end{lemma}

  \begin{proof}
  Comme $k_0^p\subset K_0^p$, on a
  \[
    k_0^p\bigl(K_0^p((T_1,\ldots,T_r))\bigr)
    =K_0^p((T_1,\ldots,T_r)).
  \]
  Si $(c_i)_{1\leqslant i\leqslant n}$ est une base de $k_0$ sur $k_0^p$,
  il est immédiat que $(c_i)$ est aussi un système de générateurs de chacun
  des deux membres de (21.9.6.8) sur $K_0^p((T_1,\ldots,T_r))$.
  \end{proof}
\end{enumerate}
\end{proof}

\begin{remark}[21.9.7]
\label{0.21.9.7-fr}
L'assertion de (21.9.6, (ii)) \emph{ne s'étend pas} au cas où
$[k_0:k_0^p]$ est infini. Prenons en effet pour $k_0$ le corps
$\mathbf F_p(X_n)_{n\geqslant0}$ des fractions rationnelles à une infinité
d'indéterminées, de sorte que les $X_n$ forment une $p$-base de $k_0$ sur
$\mathbf F_p$, et par suite $[k_0:k_0^p]=+\infty$. Dans l'énoncé de (21.9.6),
prenons $K_0=k_0$ (donc nécessairement $k=k_0$), $L_0=k_0((T))$, pour $E$
l'extension $L_0(y)$, où $y$ est racine du polynôme
\[
  Y^p-\sum_{n=0}^{\infty}X_nT^{np}
\]
de $L_0[Y]$. Alors la différence $\delta$ des deux membres de (21.9.6.2) est
\emph{non nulle}. En effet, $E$ est \emph{séparable} sur $k_0$, sans quoi
$y^p$ serait élément de $k_0(L_0^p)$, et l'on voit aussitôt que cela n'est
pas possible (dû au fait qu'il y a une infinité de $X_n$ algébriquement
indépendants)~; on a donc $d(E/K_0,k_0)=0$, la formule (21.9.6.3) se réduit à
$\delta=d(E/M_0,M_0)-1$, et il est clair que
$d(E/M_0,M_0)=\operatorname{rg}_E\Upsilon_{E/M_0}$. Tout revient à vérifier
que ce rang est $2$. Or on a $E^p\subset M_0$, une $p$-base de $E$ sur $L_0$
est formée du seul élément $y$, et $L_0$ a évidemment une $p$-base sur $M_0$
formée du seul élément $T$~; donc, puisque $E$ est extension algébrique de
$M_0$, notre assertion résulte de (21.7.1) et (21.4.5).
\end{remark}

\begin{proposition}[21.9.8]
\label{0.21.9.8-fr}
Soient $k_0$ un corps (de caractéristique quelconque), $K_0$ une extension
séparable de $k_0$, $A$ un anneau local noethérien complet, qui est une
$k_0$-algèbre et dont le corps résiduel

